\documentclass[12pt,a4paper,french]{report}

% Packages essentiels
\usepackage[french]{babel}
\usepackage[utf-8]{inputenc}
\usepackage[T1]{fontenc}
\usepackage{geometry}
\usepackage{hyperref}
\usepackage{graphicx}
\usepackage{listings}
\usepackage{xcolor}
\usepackage{fancyhdr}
\usepackage{tocloft}
\usepackage{array}
\usepackage{booktabs}
\usepackage{float}

% Configuration des marges
\geometry{
    left=2.5cm,
    right=2.5cm,
    top=2.5cm,
    bottom=2.5cm
}

% Configuration des couleurs pour le code
\definecolor{codegreen}{rgb}{0,0.6,0}
\definecolor{codegray}{rgb}{0.5,0.5,0.5}
\definecolor{codepurple}{rgb}{0.58,0,0.82}
\definecolor{backcolour}{rgb}{0.95,0.95,0.92}

% Configuration des listings pour PHP
\lstdefinelanguage{PHP}{
    language=PHP,
    basicstyle=\ttfamily\small,
    backgroundcolor=\color{backcolour},
    commentstyle=\color{codegreen},
    keywordstyle=\color{codepurple},
    numberstyle=\tiny\color{codegray},
    stringstyle=\color{codepurple},
    breakatwhitespace=false,
    breaklines=true,
    captionpos=b,
    keepspaces=true,
    numbers=left,
    numbersep=5pt,
    showspaces=false,
    showstringspaces=false,
    showtabs=false,
    tabsize=4,
    morekeywords={declare, strict_types}
}

% En-têtes et pieds de page
\pagestyle{fancy}
\fancyhf{}
\rhead{Année académique 2025-2026}
\lhead{Système de Facturation Supermarché}
\cfoot{\thepage}

% Configuration de la table des matières
\renewcommand{\cftsecleader}{\cftdotfill{\cftdotsep}}

% Titre personnalisé pour les chapitres
\usepackage{titlesec}
\titleformat{\chapter}[display]
    {\normalfont\Large\bfseries}
    {\chaptertitlename\ \thechapter}{20pt}{\Large}

\title{%
    \textbf{PAYPAS} \\[0.5cm]
    \textbf{travail pratique académique} \\[1cm]
    Développement Web en PHP Procédural \\[0.5cm]
}
\author{CODE DREAMEUR -- Faculté de Sciences Informatiques}
\date{28 avril 2026}

\begin{document}

% Page de titre
\maketitle

\newpage

% Table des matières
\tableofcontents

\newpage

% ============================================================================
\chapter{Introduction}
% ============================================================================

\section{Contexte et problématique}

Un petit supermarché souhaite moderniser son processus de vente. Actuellement, les caissiers saisissent manuellement 
les articles, ce qui génère des erreurs et ralentit les transactions. Le gérant souhaite mettre en place un système de 
caisse informatisé, accessible depuis un navigateur web, capable de lire des codes-barres via la caméra d'un téléphone 
ou d'un ordinateur, et de produire des factures détaillées.

\section{Objectifs pédagogiques}

Ce projet vise à développer les compétences suivantes :

\begin{enumerate}
    \item Concevoir et structurer une application web PHP selon le paradigme procédural;
    \item Utiliser les fichiers comme mécanisme de persistance des données;
    \item Implémenter un système d'authentification et de contrôle d'accès basé sur les rôles;
    \item Intégrer une bibliothèque de lecture de codes-barres dans une interface web;
    \item Appliquer les bonnes pratiques de validation et d'assainissement des données en PHP;
    \item Rédiger un rapport technique structuré en \LaTeX{}
\end{enumerate}

\section{Contraintes du projet}

Toute la persistance des données doit être assurée exclusivement par des fichiers au format JSON.
 L'utilisation d'un système de gestion de base de données (MySQL, PostgreSQL, SQLite, MongoDB, etc.) 
 est strictement interdite.

% ============================================================================
\chapter{Architecture et conception}
% ============================================================================

\section{Présentation générale}

L'application est développée en PHP procédural, sans utilisation de base de données. 
Elle suit une architecture en couches avec une séparation claire entre l'authentification, la logique métier et 
la présentation.

\section{Arborescence du projet}

L'organisation des fichiers est la suivante :

\begin{verbatim}
facturation/
├── index.php              # Tableau de bord principal
├── rapport.tex            # Ce rapport
├── assets/                # Ressources statiques
│   ├── css/
│   │   └── style.css      # Feuille de styles
│   ├── images/            # Images du projet
│   └── js/
│       └── scanner.js     # Bibliothèque de lecture codes-barres
├── auth/                  # Gestion authentification
│   ├── login.php
│   ├── logout.php
│   └── session.php
├── config/                # Configuration globale
│   └── config.php
├── data/                  # Fichiers de persistance
│   ├── factures.json
│   ├── produits.json
│   ├── users.json
│   └── utilisateurs.json
├── includes/              # Fonctions réutilisables
│   ├── fonctions-auth.php
│   ├── fonctions-factures.php
│   ├── fonctions-produits.php
│   ├── header.php
│   ├── footer.php
│   └── helpers.php
├── modules/               # Modules fonctionnels
│   ├── caisse.php
│   ├── factures.php
│   ├── inscription.php
│   ├── produits.php
│   ├── admin/             # Fonctionnalités admin
│   │   ├── ajouter-compte.php
│   │   ├── gestion-comptes.php
│   │   └── supprimer-comptes.php
│   ├── facturation/       # Factures
│   │   ├── afficher-facture.php
│   │   ├── calcul.php
│   │   └── nouvelle-facture.php
│   ├── produits/          # Gestion produits
│   │   ├── enregistrer.php
│   │   ├── lire.php
│   │   └── liste.php
└── rapports/              # Génération rapports
    ├── rapport-journalier.php
    └── rapport-mensuel.php
\end{verbatim}

\section{Choix technologiques}

\subsection{PHP procédural}
Le choix du paradigme procédural permet une compréhension progressive de la programmation web sans abstraction excessive. Le code reste lisible et direct.

\subsection{Fichiers JSON pour la persistance}
Les fichiers JSON offrent une structure simple et lisible, facilement éditable et sans dépendance externe.
 Chaque entité (produits, factures, utilisateurs) est stockée dans un fichier JSON séparé.

\subsection{Bibliothèque Html5Qrcode}
Cette bibliothèque JavaScript permet la lecture des codes-barres directement via la caméra du navigateur,
 sans installation de plugin.

% ============================================================================
\chapter{Réalisation des objectifs}
% ============================================================================

\section{Objectif 1 : Architecture PHP procédurale}

L'application est entièrement structurée selon le paradigme procédural :

\begin{itemize}
    \item \textbf{Configuration centralisée} (\texttt{config/config.php}) : constantes globales, chemins des fichiers
    \item \textbf{Fonctions utilitaires} (\texttt{includes/helpers.php}) : fonctions réutilisables
    \item \textbf{Séparation des responsabilités} : authentification, métier, présentation
    \item \textbf{Modules indépendants} : chaque module gère un domaine spécifique
\end{itemize}

\subsection{Exemple : Helpers}

\begin{lstlisting}[language=PHP, caption=Fonctions essentielles (helpers.php)]
function e(string $value): string
{
    return htmlspecialchars($value, ENT_QUOTES, 'UTF-8');
}

function read_json_file(string $filePath): array
{
    if (!file_exists($filePath)) return [];
    $content = file_get_contents($filePath);
    $decoded = json_decode($content, true);
    return is_array($decoded) ? $decoded : [];
}

function write_json_file(string $filePath, array $data): bool
{
    $json = json_encode($data, JSON_PRETTY_PRINT | JSON_UNESCAPED_UNICODE);
    return file_put_contents($filePath, $json) !== false;
}
end{lstlisting}

\section{Objectif 2 : Persistance par fichiers}

Les données sont entièrement persistées en JSON :

\subsection{Structure des fichiers}

\subsubsection{Produits (produits.json)}
\begin{lstlisting}[language=PHP]
[
    {
        "code": "100001",
        "nom": "Produit A",
        "prix": 15.99,
        "stock": 100
    }
]
\end{lstlisting}

\subsubsection{Factures (factures.json)}
\begin{lstlisting}[language=PHP]
[
    {
        "numero": "FAC-20260427-0001",
        "date": "2026-04-27 12:27:01",
        "caissier": "admin",
        "lignes": [
            {
                "code": "100001",
                "nom": "Produit A",
                "prix": 15.99,
                "quantite": 2
            }
        ],
        "total": 31.98
    }
]
\end{lstlisting}

\subsubsection{Utilisateurs (users.json)}
\begin{lstlisting}[language=PHP]
[
    {
        "username": "admin",
        "password": "admin123",
        "nom": "Administrateur",
        "role": "admin"
    },
    {
        "username": "caissier",
        "password": "caissier123",
        "nom": "Caissier 1",
        "role": "caissier"
    }
]
\end{lstlisting}

\subsection{Opérations de persistance}

\begin{lstlisting}[language=PHP, caption=Lecture et écriture JSON]
// Lecture
$produits = read_json_file(PRODUCTS_FILE);

// Écriture
$factures = read_json_file(INVOICES_FILE);
$factures[] = $nouvelle_facture;
write_json_file(INVOICES_FILE, $factures);
\end{lstlisting}

\section{Objectif 3 : Authentification et contrôle d'accès par rôles}

\subsection{Système de sessions}

L'authentification utilise les sessions PHP :

\begin{lstlisting}[language=PHP, caption=Gestion des sessions (auth/session.php)]
session_start();

function is_logged_in(): bool
{
    return isset($_SESSION['user']) && is_array($_SESSION['user']);
}

function current_user(): ?array
{
    return $_SESSION['user'] ?? null;
}

function require_login(): void
{
    if (!is_logged_in()) {
        redirect_to('auth/login.php');
    }
}
\end{lstlisting}

\subsection{Authentification utilisateur}

Les identifiants sont vérifiés dans \texttt{auth/login.php} :

\begin{lstlisting}[language=PHP, caption=Authentification (extrait)]
$users = read_json_file(USERS_FILE);

foreach ($users as $user) {
    if ($user['username'] === $username && 
        $user['password'] === $password) {
        $_SESSION['user'] = [
            'username' => $user['username'],
            'nom' => $user['nom'],
            'role' => $user['role'],
        ];
        redirect_to('index.php');
    }
}
\end{lstlisting}

\subsection{Contrôle d'accès par rôles}

Les rôles sont stockés dans la session et peuvent être vérifiés :

\begin{itemize}
    \item \textbf{admin} : accès complet, gestion des comptes
    \item \textbf{caissier} : accès à la caisse, lecture codes-barres
    \item \textbf{gérant} : accès aux rapports et statistiques
\end{itemize}

\section{Objectif 4 : Lecture de codes-barres}

\subsection{Intégration Html5Qrcode}

La bibliothèque est intégrée dans \texttt{assets/js/scanner.js} :

\begin{lstlisting}[language=PHP, caption=Scanner de codes-barres]
document.addEventListener("DOMContentLoaded", function () {
    const scanBtn = document.getElementById("start-scan");
    const codeInput = document.getElementById("code-input");

    scanBtn.addEventListener("click", async function () {
        let scanner = new Html5Qrcode("scanner-video");
        
        await scanner.start(
            { facingMode: "environment" },
            { fps: 10, qrbox: 250 },
            function (decodedText) {
                codeInput.value = decodedText;
                scanner.stop();
            }
        );
    });
});
\end{lstlisting}

\subsection{Fonctionnalités}

\begin{itemize}
    \item Activation par clic sur le bouton « Scanner »
    \item Utilisation de la caméra de l'appareil
    \item Décodage automatique du code-barres
    \item Remplissage automatique du champ de saisie
    \item Fallback : saisie manuelle possible
\end{itemize}

\section{Objectif 5 : Validation et assainissement des données}

\subsection{Fonctions de validation}

Les données POST sont validées via des fonctions typées :

\begin{lstlisting}[language=PHP, caption=Fonctions de validation (helpers.php)]
function post_string(string $key): string
{
    $value = $_POST[$key] ?? '';
    return trim((string) $value);
}

function post_float(string $key): float
{
    $value = str_replace(',', '.', post_string($key));
    return is_numeric($value) ? (float) $value : 0.0;
}

function post_int(string $key): int
{
    return (int) post_string($key);
}
\end{lstlisting}

\subsection{Assainissement en sortie}

Les données en sortie HTML sont échappées :

\begin{lstlisting}[language=PHP]
function e(string $value): string
{
    return htmlspecialchars($value, ENT_QUOTES, 'UTF-8');
}

// Utilisation en template
<p><?= e($user_input); ?></p>
\end{lstlisting}

\subsection{Exemples d'application}

\begin{lstlisting}[language=PHP, caption=Validation lors de l'ajout de produit]
if ($_SERVER['REQUEST_METHOD'] === 'POST') {
    $code = post_string('code');           // Chaîne validée
    $nom = post_string('nom');             // Chaîne validée
    $prix = post_float('prix');            // Float validé
    $stock = post_int('stock');            // Entier validé
    
    // Vérification logique
    if ($prix < 0) {
        $error = 'Le prix ne peut pas être négatif';
    }
}
\end{lstlisting}

\subsection{Déclaration de types stricts}

Chaque fichier PHP utilise la déclaration de types stricts :

\begin{lstlisting}[language=PHP]
<?php
declare(strict_types=1);

// Évite les conversions de type implicites
\end{lstlisting}

% ============================================================================
\chapter{Implémentation détaillée}
% ============================================================================

\section{Modules de caisse}

\subsection{Module caisse.php}

La caisse permet d'ajouter des produits au panier et de générer des factures.

\begin{itemize}
    \item \textbf{Saisie de codes-barres} : Scanner ou saisie manuelle
    \item \textbf{Ajout au panier} : Vérification du stock
    \item \textbf{Gestion du panier} : Modification, suppression, vider
    \item \textbf{Validation} : Avant chaque action
    \item \textbf{Génération de facture} : Création et sauvegarde JSON
\end{itemize}

\subsection{Processus de vente}

\begin{enumerate}
    \item Caissier se connecte
    \item Caissier scanne ou saisit un code-barres
    \item Produit ajouté au panier (vérification stock)
    \item Répétition des étapes 2-3
    \item Clôture de la facture
    \item Facture sauvegardée, stock mis à jour
\end{enumerate}

\section{Gestion des produits}

\subsection{Opérations disponibles}

\begin{itemize}
    \item \textbf{Créer} : Ajouter un nouveau produit avec code-barres
    \item \textbf{Lire} : Afficher la liste des produits et stocks
    \item \textbf{Mettre à jour} : Modifier prix, stock, nom
    \item \textbf{Supprimer} : Retirer un produit du catalogue
\end{itemize}

\subsection{Format des données}

Chaque produit contient :

\begin{table}[H]
\centering
\begin{tabular}{|l|l|l|}
\hline
\textbf{Champ} & \textbf{Type} & \textbf{Description} \\
\hline
code & string & Code-barres unique \\
nom & string & Nom du produit \\
prix & float & Prix unitaire (euros) \\
stock & integer & Quantité en stock \\
\hline
\end{tabular}
\caption{Structure d'un produit}
\end{table}

\section{Gestion des factures}

\subsection{Génération de factures}

Les factures sont générées avec un numéro unique au format : \texttt{FAC-YYYYMMDD-XXXX}

Exemple : \texttt{FAC-20260427-0001}

\subsection{Contenu d'une facture}

\begin{table}[H]
\centering
\begin{tabular}{|l|l|}
\hline
\textbf{Champ} & \textbf{Description} \\
\hline
numero & Identifiant unique \\
date & Timestamp de création \\
caissier & Nom de l'utilisateur \\
lignes & Tableau des articles \\
total & Montant total (euros) \\
\hline
\end{tabular}
\caption{Structure d'une facture}
\end{table}

\section{Administration}

\subsection{Gestion des comptes}

Les administrateurs peuvent :

\begin{itemize}
    \item Ajouter de nouveaux utilisateurs
    \item Modifier les informations utilisateur
    \item Supprimer des comptes
    \item Gérer les rôles
\end{itemize}

\subsection{Utilisateurs par défaut}

\begin{table}[H]
\centering
\begin{tabular}{|l|l|l|}
\hline
\textbf{Username} & \textbf{Password} & \textbf{Rôle} \\
\hline
admin & admin & Admin \\
caissier & caissier & Caissier \\
\hline
\end{tabular}
\caption{Comptes de démonstration}
\end{table}

% ============================================================================
\chapter{Tests et résultats}
% ============================================================================

\section{Scénarios de test}

\subsection{Test 1 : Authentification}

\textbf{Scénario :} Un utilisateur se connecte avec des identifiants valides

\textbf{Résultat :} L'utilisateur est redirigé vers le tableau de bord, la session est créée

\subsection{Test 2 : Ajout de produit}

\textbf{Scénario :} Ajout d'un produit via le formulaire

\textbf{Résultat :} Le produit est enregistré dans produits.json, le stock est défini

\subsection{Test 3 : Lecture code-barres}

\textbf{Scénario :} Activation du scanner et scannage d'un code-barres

\textbf{Résultat :} Le code est décodé et inséré dans le champ

\subsection{Test 4 : Création facture}

\textbf{Scénario :} Ajout de produits au panier et clôture

\textbf{Résultat :} Facture créée avec numéro unique, stock mis à jour

\subsection{Test 5 : Validation des données}

\textbf{Scénario :} Saisie de données invalides (prix négatif, stock non numérique)

\textbf{Résultat :} Données rejetées, message d'erreur affichée

\subsection{Test 6 : Persistance}

\textbf{Scénario :} Rechargement de la page, redémarrage du serveur

\textbf{Résultat :} Les données persisten t via les fichiers JSON

\section{Limitations identifiées}

\begin{enumerate}
    \item Les mots de passe ne sont pas chiffrés (démo uniquement)
    \item Pas de mécanisme de verrouillage après tentatives échouées
    \item Pas de gestion des transactions multi-utilisateurs
    \item Pas de sauvegarde/archivage automatique
    \item La lecture de codes-barres nécessite une caméra fonctionnelle
\end{enumerate}

% ============================================================================
\chapter{Conclusion}
% ============================================================================

\section{Bilan du projet}

Ce projet a permis de développer une application web fonctionnelle en PHP procédural, respectant les contraintes imposées :

\begin{itemize}
    \item ✓ Architecture procédurale clara et organisée
    \item ✓ Persistance par fichiers JSON
    \item ✓ Authentification avec rôles
    \item ✓ Lecture de codes-barres intégrée
    \item ✓ Validation robuste des données
    \item ✓ Rapport technique en \LaTeX{}
\end{itemize}

\section{Compétences acquises}

\begin{enumerate}
    \item Structuration d'une application web procédurale
    \item Gestion des sessions et authentification
    \item Manipulation de fichiers et format JSON
    \item Validation et sécurité des données
    \item Intégration de bibliothèques externes (JavaScript)
    \item Rédaction de documentation technique
\end{enumerate}

\section{Améliorations futures}

\subsection{Évolutions technologiques}

\begin{itemize}
    \item Migration vers une base de données (PostgreSQL)
    \item Architecture orientée objet (POO)
    \item Framework comme Symfony ou Laravel
    \item API REST pour applications mobiles
\end{itemize}

\subsection{Fonctionnalités additionnelles}

\begin{itemize}
    \item Système de remise/promotion
    \item Gestion des clients et fidélité
    \item Intégration de moyens de paiement
    \item Synchronisation multi-caisses
    \item Export comptable (fichiers SAE)
    \item Alertes de stock bas
    \item Historique d'audit complet
\end{itemize}

% ============================================================================
% Annexes
% ============================================================================

\appendix

\chapter{Code source : Helpers}

\lstinputlisting[language=PHP, caption=Fonctions utilitaires complètes]{../includes/helpers.php}

\chapter{Configuration}

\begin{lstlisting}[language=PHP, caption=Fichier de configuration]
<?php
declare(strict_types=1);

define('USERS_FILE', __DIR__ . '/../data/users.json');
define('PRODUCTS_FILE', __DIR__ . '/../data/produits.json');
define('INVOICES_FILE', __DIR__ . '/../data/factures.json');
define('TAX_RATE', 0.20); // TVA 20%
define('DATE_FORMAT', 'Y-m-d H:i:s');
\end{lstlisting}

\chapter{Captures d'écran}

\section{Tableau de bord}

Le tableau de bord affiche les modules disponibles selon le rôle de l'utilisateur :
\begin{itemize}
    \item Scanner/Caisse
    \item Ajouter des produits
    \item Consulter les factures
    \item Ajouter un compte (admin)
\end{itemize}

\section{Interface de caisse}

L'interface affiche :
\begin{itemize}
    \item Champ de saisie/scan du code-barres
    \item Bouton d'activation du scanner
    \item Panier avec articles ajoutés
    \item Total à payer
    \item Bouton de clôture
\end{itemize}

\chapter{Ressources et références}

\section{Documentation PHP}

\begin{itemize}
    \item Sessions : \url{https://www.php.net/manual/fr/book.session.php}
    \item Fichiers : \url{https://www.php.net/manual/fr/book.filesystem.php}
    \item JSON : \url{https://www.php.net/manual/fr/book.json.php}
\end{itemize}

\section{Bibliothèques externes}

\begin{itemize}
    \item Html5Qrcode : \url{https://github.com/mebjas/html5-qrcode}
\end{itemize}

\end{document}
